% ---
% id: meat-core/lem:cone-concavity
% kind: lem
% title: "Concavity on the positive cone"
% summary: >-
%   The pinched trace-power functional T is positively homogeneous and concave
%   on the whole positive cone. Hence the sector deficits are concave on the
%   cone and the branch quantities K on its sector-block-diagonal part, for
%   every block-diagonal positive measurement family, complete or not.
% prominence: standard
% status: accepted
% uses: [thm:branch-formula, foundations/fact:concavity, cor:concavity, lem:branch-maps, def:maps, def:admissible-declaration]
% mentions: [meat-core/frank-wolfe/thm:firstorder, meat-core/frank-wolfe/prop:relaxed-set-properties]
% assumes: []
% refs: []
% provenance: newly-derived
% lean:
% notes: "Parts (a) to (c) are cited by letter from other claims of this contract; keep the lettering."
% ---

Corollary~\ref{cor:concavity} gives concavity only on subnormalised
operators $\Sleq{}$, a domain restriction it inherits from
Lemma~\ref{lem:branch-maps}\,(i). Positive homogeneity makes the restriction
dispensable: two positive operators, divided by the larger of their
traces, lie in $\Sleq{}$, and $T$ scales back linearly. Dispensing with it
is what part~(c) needs, and what every use of $T$ off the trace-one slice
needs, the relaxed iterate set of the Frank--Wolfe submodule
(Proposition~\ref{prop:relaxed-set-properties}\,(c)) among them. Part~(c) reads no
completeness relation of $\{F_v\}$, which is what makes it available off
that slice as well. This lemma is the only place where the positive
homogeneity of $T$ is established, and every later use of it, the
objective of Theorem~\ref{thm:firstorder} included, reads it from here.

\begin{lemma}\label{lem:cone-concavity}
Let $\alpha\in(1,\infty)$, put $p\coloneqq1/\alpha\in(0,1)$, and let $\cZ$
be a pinching built from a complete family of mutually orthogonal
projectors. On the cone $\Pos{}$ of the space $\cZ$ acts on, consider the
functional
\begin{equation}\label{eq:Tcone}
  T(X)\coloneqq\tr\!\big[\big(\cZ(X^{p})\big)^{\alpha}\big].
\end{equation}
\begin{enumerate}[label=(\alph*),leftmargin=2.4em]
\item $T$ is positively $1$-homogeneous: $T(tX)=t\,T(X)$ for every
$t\ge0$.
\item $T$ is concave on all of $\Pos{}$.
\item Consider a declaration whose components other than $\{F_v\}$
satisfy Definition~\ref{def:admissible-declaration} and whose $\{F_v\}$ is any family of
sector-block-diagonal positive semidefinite operators, complete or not.
Then each deficit $f_{\hat c,s}$ of \eqref{eq:f-def} is concave on all of
$\Pos{A\Qsys_s}$, and $K_{\hat c}$ of \eqref{eq:fK} is concave on the
sector-block-diagonal part of $\Pos{A\Qsys}$, in particular on $\Feas$.
\end{enumerate}
\end{lemma}

\begin{proof}
(a) Let $t>0$. First, $(tX)^{p}=t^{p}X^{p}$: the map $X\mapsto X^{p}$ is
defined by the functional calculus of $x\mapsto x^{p}$, and
$(t\lambda)^{p}=t^{p}\lambda^{p}$ on each eigenvalue, the eigenprojectors
being unchanged by the scaling. Second, $\cZ$ is a pinching, hence
linear, so $\cZ\big((tX)^{p}\big)=\cZ(t^{p}X^{p})=t^{p}\cZ(X^{p})$, and
$\tr[(t^{p}Y)^{\alpha}]=t^{p\alpha}\tr[Y^{\alpha}]$ by the same functional
calculus applied to $x\mapsto x^{\alpha}$. Together,
$T(tX)=t^{p\alpha}T(X)=t\,T(X)$ because $p\alpha=1$. At $t=0$ both sides
vanish.

(b) Let $X,Y\in\Pos{}$ and $\lambda\in(0,1)$. If $\tr X=\tr Y=0$ then
$X=Y=0$, a positive semidefinite operator of zero trace being zero, and
the inequality reads $0\ge0$. Otherwise put
$c\coloneqq\max\{\tr X,\tr Y\}>0$, so that $X/c$ and $Y/c$ lie in
$\Sleq{}$ and so does their convex combination. \xref{foundations/fact:concavity}
applies at $\alpha'=p\in(0,1]$, the substitution
Corollary~\ref{cor:concavity} makes, and on the source's own stated
domain $\Sleq{}$, so no hypothesis of the import is stretched; it gives
\begin{equation*}
  T\Big(\lambda\frac{X}{c}+(1-\lambda)\frac{Y}{c}\Big)
  \;\ge\;\lambda\,T\Big(\frac{X}{c}\Big)
  +(1-\lambda)\,T\Big(\frac{Y}{c}\Big).
\end{equation*}
Multiply by $c>0$ and apply (a) once to each of the three arguments:
$T(\lambda X+(1-\lambda)Y)\ge\lambda T(X)+(1-\lambda)T(Y)$. No trace
hypothesis on $X$ or $Y$ is used: the scaling is by their own traces.

(c) $\cG_{\hat c,s}$ is completely positive by Lemma~\ref{lem:branch-maps}\,(i) in
its Kraus-form clause, which needs $F^{(s)}_v\succeq0$ and no completeness
relation; so it is linear and maps $\Pos{A\Qsys_s}$ into $\Pos{}$.
Precomposition of the concave $T$ of (b) with a linear map preserves
concavity wherever the image lies in $\Pos{}$, which is on all of
$\Pos{A\Qsys_s}$, and subtracting the linear
$\sigma\mapsto\tr\cG_{\hat c,s}(\sigma)$ preserves it; by
\eqref{eq:f-def} that is concavity of $f_{\hat c,s}$. For $K_{\hat c}$:
on sector-block-diagonal $\omega$, Lemma~\ref{lem:branch-maps}\,(iii) gives
$\cG_{\hat c}(\omega)=\bigoplus_s\cG_{\hat c,s}(\omega_s)$ with mutually
orthogonal summands, and powers and $\cZ$ act blockwise on that
decomposition, so $T_{\hat c}(\omega)=\sum_sT_{\hat c,s}(\omega_s)$ and
$\tr\cG_{\hat c}(\omega)=\sum_s\tr\cG_{\hat c,s}(\omega_s)$. Inserting
both into \eqref{eq:fK} gives
$K_{\hat c}(\omega)=P_{\hat c}(\omega)+\sum_sf_{\hat c,s}(\omega_s)$,
which is \eqref{eq:c5-sector} obtained from the definitions and
Lemma~\ref{lem:branch-maps}\,(iii) alone, without Theorem~\ref{thm:branch-formula}'s entropy
statement. $P_{\hat c}$ is affine by \eqref{eq:Phimap}, and a sum of
concave functions is concave.
\end{proof}
